% =======================================================
% ШАБЛОНЫ и RAII
% =======================================================

\section{Шаблоны функций и классов. Специализация}

% -------------------------------------------------------
\subsection{Зачем нужны шаблоны}

\textbf{Шаблоны} --- механизм \textbf{обобщённого программирования}: написать код
один раз для \textbf{любого типа}, а компилятор сгенерирует конкретные версии.
Это статический полиморфизм, без runtime-затрат.

\subsection{Шаблоны функций}

\begin{lstlisting}
template<typename T>
T max(T a, T b) {
    return (a > b) ? a : b;
}
max(3, 5);        // T = int    (выведен из аргументов)
max(2.5, 1.5);    // T = double
max<int>(3, 5);   // T задан явно
\end{lstlisting}

\subsubsection{Вывод аргументов шаблона (template argument deduction)}

Компилятор сам определяет \texttt{T} по типам аргументов. Если вывод
неоднозначен, тип указывают явно:

\begin{lstlisting}
template<typename T> void f(T x);
f(42);          // T = int
f<double>(42);  // T = double явно
\end{lstlisting}

\subsection{Шаблоны классов}

\begin{lstlisting}
template<typename T>
class Stack {
    std::vector<T> data_;
public:
    void push(const T& x) { data_.push_back(x); }
    T pop() { T t = data_.back(); data_.pop_back(); return t; }
    bool empty() const { return data_.empty(); }
};
Stack<int> si;          // стек целых
Stack<std::string> ss;  // стек строк
\end{lstlisting}

\subsubsection{Class Template Argument Deduction (CTAD, C++17)}

С C++17 для шаблонов классов аргументы тоже могут выводиться:

\begin{lstlisting}
std::vector v = {1, 2, 3};   // std::vector<int>, тип выведен
std::pair  p{1, 2.0};        // std::pair<int, double>
\end{lstlisting}

\subsection{Параметры шаблона}

\begin{itemize}[leftmargin=2em]
    \item \textbf{Типовые:} \texttt{template<typename T>}.
    \item \textbf{Нетиповые (NTTP):} значения, известные на этапе компиляции:
          \texttt{template<int N>}, \texttt{template<auto V>}.
    \item \textbf{Шаблонные:} шаблон как параметр шаблона.
\end{itemize}

\begin{lstlisting}
template<typename T, size_t N>
class Array {
    T data_[N];   // N -- константа времени компиляции
public:
    constexpr size_t size() const { return N; }
};
Array<int, 5> a;   // массив на 5 int
\end{lstlisting}

\subsection{Полная специализация}

Особая реализация шаблона для \textbf{конкретного} типа:

\begin{lstlisting}
template<typename T>
struct TypeName { static const char* get() { return "unknown"; } };

template<>   // полная специализация для bool
struct TypeName<bool> { static const char* get() { return "bool"; } };

template<>   // полная специализация для int
struct TypeName<int> { static const char* get() { return "int"; } };
\end{lstlisting}

\subsection{Частичная специализация}

Специализация не для одного типа, а для \textbf{семейства} типов. Доступна
\textbf{только для шаблонов классов} (не функций):

\begin{lstlisting}
template<typename T>
struct IsPointer { static constexpr bool value = false; };

template<typename T>            // частичная: для любых указателей T*
struct IsPointer<T*> { static constexpr bool value = true; };

IsPointer<int>::value;    // false
IsPointer<int*>::value;   // true
\end{lstlisting}

\textbf{Для функций} вместо частичной специализации используют перегрузку.

\subsection{Двухфазная компиляция и \texttt{typename}/\texttt{template}}

Шаблон компилируется в два прохода: сначала проверяется код, не зависящий от
параметров, при инстанцировании --- зависимый. Перед зависимым именем типа нужен
\texttt{typename}:

\begin{lstlisting}
template<typename T>
void f() {
    typename T::value_type x;   // typename: говорим "это тип"
    typename std::vector<T>::iterator it;
}
\end{lstlisting}

\subsection{Где определять шаблоны}

Определение шаблона должно быть видимо в точке инстанцирования, поэтому шаблоны
обычно пишут \textbf{целиком в заголовках} (а не разделяют на \texttt{.h}/\texttt{.cpp}).

\subsection{Переменные-шаблоны и \texttt{constexpr}}

\begin{lstlisting}
template<typename T>
constexpr T pi = T(3.1415926535897932385);
auto x = pi<float>;   // 3.14159f
\end{lstlisting}

\subsection{Концепты (C++20) --- кратко}

Концепты ограничивают допустимые типы и дают понятные ошибки (подробно --- в
разделе о метапрограммировании):

\begin{lstlisting}
#include <concepts>
template<std::integral T>   // только целочисленные T
T half(T x) { return x / 2; }
\end{lstlisting}


% =======================================================
\section{RAII. Умные указатели}

% -------------------------------------------------------
\subsection{Идиома RAII}

\textbf{RAII (Resource Acquisition Is Initialization)} --- ключевая идиома C++:
\textbf{ресурс захватывается в конструкторе и освобождается в деструкторе}.
Поскольку деструктор вызывается автоматически (в т.\,ч. при исключении), ресурс
не утечёт.

\begin{lstlisting}
class FileGuard {
    FILE* f_;
public:
    explicit FileGuard(const char* name) : f_(fopen(name, "r")) {}
    ~FileGuard() { if (f_) fclose(f_); }   // освобождение гарантировано
    FileGuard(const FileGuard&) = delete;  // не копируем владение
    FILE* get() const { return f_; }
};

void process() {
    FileGuard fg("data.txt");
    // ... даже если здесь бросится исключение ...
}   // ~FileGuard() вызовется и закроет файл в любом случае
\end{lstlisting}

\textbf{Под RAII попадают:} память, файлы, мьютексы (\texttt{std::lock\_guard}),
сокеты, соединения с БД --- любой ресурс с парой «захват/освобождение».

\subsection{Проблема сырых указателей}

\begin{lstlisting}
void leak() {
    Widget* w = new Widget();
    if (error()) return;   // УТЕЧКА: delete не вызовется
    delete w;
}
\end{lstlisting}

Умные указатели --- это RAII для динамической памяти: они \textbf{сами} вызывают
\texttt{delete}.

\subsection{\texttt{std::unique\_ptr} --- единоличное владение}

Владеет объектом эксклюзивно. Нельзя копировать (только перемещать). Нулевые
накладные расходы по сравнению с сырым указателем.

\begin{lstlisting}
#include <memory>
std::unique_ptr<Widget> p = std::make_unique<Widget>(args);
p->method();              // как обычный указатель
// auto q = p;            // ОШИБКА: копирование запрещено
auto q = std::move(p);    // передача владения: теперь p == nullptr
// при выходе из области видимости q сам сделает delete
\end{lstlisting}

\textbf{Используйте \texttt{make\_unique}}, а не \texttt{new}: короче, безопаснее
при исключениях, без дублирования типа.

\subsubsection{unique\_ptr для массивов и кастомный делитер}

\begin{lstlisting}
std::unique_ptr<int[]> arr = std::make_unique<int[]>(10);  // массив
std::unique_ptr<FILE, decltype(&fclose)> f(fopen("x","r"), &fclose);
\end{lstlisting}

\subsection{\texttt{std::shared\_ptr} --- разделяемое владение}

Несколько \texttt{shared\_ptr} владеют одним объектом совместно. Внутри ---
\textbf{счётчик ссылок (reference count)}. Объект уничтожается, когда счётчик
доходит до нуля.

\begin{lstlisting}
std::shared_ptr<Widget> p1 = std::make_shared<Widget>();
std::shared_ptr<Widget> p2 = p1;   // счётчик = 2 (копирование разрешено)
std::cout << p1.use_count();       // 2
p2.reset();                        // счётчик = 1
// объект жив, пока есть хоть один владелец
\end{lstlisting}

\textbf{\texttt{make\_shared} эффективнее}: выделяет объект и управляющий блок
(со счётчиком) одним аллокатором.

\subsubsection{Накладные расходы}

\texttt{shared\_ptr} весит как два указателя (на объект и на управляющий блок),
а изменение счётчика --- \textbf{атомарная} операция (потокобезопасно, но не
бесплатно). Если разделение не нужно --- используйте \texttt{unique\_ptr}.

\subsection{\texttt{std::weak\_ptr} --- слабая ссылка}

\texttt{weak\_ptr} наблюдает за объектом \texttt{shared\_ptr}, \textbf{не увеличивая
счётчик}. Чтобы получить доступ, его «повышают» через \texttt{lock()}.

\begin{lstlisting}
std::shared_ptr<Widget> sp = std::make_shared<Widget>();
std::weak_ptr<Widget> wp = sp;     // не владеет

if (auto locked = wp.lock()) {     // shared_ptr, если объект ещё жив
    locked->method();
} else {
    // объект уже уничтожен
}
\end{lstlisting}

\subsubsection{Решение проблемы циклических ссылок}

Два \texttt{shared\_ptr}, ссылающиеся друг на друга, создают \textbf{цикл} ---
счётчики никогда не дойдут до нуля, память утечёт. Один из указателей делают
\texttt{weak\_ptr}:

\begin{lstlisting}
struct Node {
    std::shared_ptr<Node> next;   // владеет следующим
    std::weak_ptr<Node>   prev;   // НЕ владеет предыдущим -> нет цикла
};
\end{lstlisting}

\subsection{Сравнение умных указателей}

\begin{center}
\begin{tabular}{lp{4cm}p{4.5cm}}
\toprule
& \textbf{Владение} & \textbf{Когда использовать} \\
\midrule
\texttt{unique\_ptr} & единоличное & по умолчанию; один владелец \\
\addlinespace
\texttt{shared\_ptr} & разделяемое & несколько владельцев \\
\addlinespace
\texttt{weak\_ptr}   & не владеет  & наблюдение, разрыв циклов \\
\bottomrule
\end{tabular}
\end{center}

\subsection{Передача умных указателей в функции}

\begin{itemize}[leftmargin=2em]
    \item Функция \textbf{забирает владение} --- \texttt{unique\_ptr<T>} по
          значению (передавать через \texttt{std::move}).
    \item Функция \textbf{разделяет владение} --- \texttt{shared\_ptr<T>} по
          значению.
    \item Функция \textbf{только использует} объект --- передавайте
          \texttt{T*} или \texttt{T\&}, а не умный указатель.
\end{itemize}

\begin{lstlisting}
void use(const Widget& w);            // просто использует -- ссылка
void take(std::unique_ptr<Widget> w); // забирает владение
take(std::move(p));
\end{lstlisting}
